\chapter{Hodnotiaci modul: ako nástroj overuje a vyhodnocuje žiakov}
\label{chapter:hodnotiaci-modul}

Táto príloha opisuje tú časť riešenia, ktorá zo \emph{návrhu merania} robí
\emph{vykonateľné meranie}. Metodika (kapitola~\ref{chapter:metodika}) hovorí,
\emph{čo} sa meria a \emph{prečo}; táto príloha hovorí, čím a ako sa to
administruje, skóruje, kontroluje a komu sa vracia. Kód je v priečinku
\texttt{assessment/}, používa výhradne štandardnú knižnicu Pythonu a je overený
\instrTests{}~automatickými testami v~\instrTestSections{}~oddieloch.

\section{Východisková zásada: položka je údaj, nie próza}
\label{sec:bank-zasada}

Pred vznikom tohto modulu existoval každý merací nástroj práve raz~--- ako text
prílohy~\ref{chapter:instrumenty}. To stačí na \emph{zdokumentovanie} testu a
nestačí na jeho \emph{administrovanie}, takže by musela vzniknúť druhá kópia: tá,
ktorú žiak skutočne vidí. Dve kópie jednej špecifikácie sú v tejto práci už dvakrát
zaplatenou chybou~--- taxonómia PEQI v~Pythone a v~Jave sa ticho rozišla a oprava
šla len do jednej; úprava kódovania konzoly šla do jedného vstupného bodu z~osemnástich.
Druhá kópia vždy prehrá, prehrá neviditeľne, a tu by prehrávajúcou kópiou bola
práve tá, podľa ktorej sa práca hodnotí.

Položka je preto údaj. Jeden zásobník, tri odberatelia, a ani jeden z~nich nedrží
vlastnú kópiu:

\begin{center}
\small
\begin{tabular}{@{}ll@{}}
\toprule
\texttt{assessment/banks/*.json} & \textbf{jediný zdroj pravdy} \\
\midrule
$\rightarrow$ \texttt{runner.py}     & to, čo žiak vidí na obrazovke \\
$\rightarrow$ \texttt{scoring.py}    & to, čo z~odpovede robí skóre \\
$\rightarrow$ \texttt{tex\_export.py} & to, čo vytlačí príloha~\ref{chapter:instrumenty} \\
\bottomrule
\end{tabular}
\end{center}

Zásobník obsahuje \asmBanks{}~nástrojov a \asmItems{}~položiek. Tretí odberateľ je
overovaný oproti druhému: každé zadanie, ktoré príloha tlačí, musí byť zadanie,
ktoré zásobník drží, a preklad práce zlyhá, ak to tak nie je. Príloha sa od
nástroja rozísť nemôže, pretože príloha z~nástroja \emph{vzniká}.

\paragraph{Čo zásobník odmieta.} Validátor nie je štylistický. Každá kontrola
zodpovedá spôsobu, akým školská štúdia stráca údaje nenávratne:

\begin{itemize}[nosep]
    \item položka bez kľúča~--- do analýzy by prišla ako stĺpec nul, čo sa nedá
          odlíšiť od veľmi ťažkej položky;
    \item rozptyľovač bez diagnózy~--- nesprávna odpoveď, z~ktorej sa nič
          nedozvieme, hoci je to jediná príležitosť sa žiaka na niečo spýtať;
    \item jazdec sebadôvery na položke, ktorá nemá správnu odpoveď;
    \item forma, ktorá nesedí na špecifikačnú tabuľku, alebo má iný počet
          rubrikových položiek než druhá forma;
    \item nástroj, ktorého položky nevojdú do svojho času v~protokole;
    \item kotviaca položka, ktorá nie je v~oboch formách znak po znaku totožná.
\end{itemize}

Posledná z~nich je dôvod, pre ktorý tento modul vznikol, a je predmetom
časti~\ref{sec:kotvenie-oprava}.

\section{Identifikácia žiaka na piatich príležitostiach}
\label{sec:kod-ziaka}

Primárny výstup štúdie je \emph{zmena}: čo žiak dosiahol v~šiestom týždni mínus
to, čo dosiahol v~nultom. Vypočítať ju znamená spojiť päť príležitostí jednej
osoby. Chrániť žiaka znamená, že výskumník tú osobu nesmie vedieť pomenovať.
Protokol (príloha~\ref{chapter:protokol}) rieši druhú polovicu dobre~--- kontaktná
osoba školy rozdá obálky s~kódmi a je jediná, kto vie priradiť kód k~žiakovi~---
a prvú polovicu necháva na prsty žiaka.

Tam bola medzera. Kód sa pri 330~žiakoch a piatich príležitostiach píše asi
1\,650-krát, pätnásťročnými, na školských klávesniciach, pod časovým tlakom.
Preklep nevyrobí chybu. Vyrobí dve:

\begin{itemize}[nosep]
    \item jeden žiak tú príležitosť ticho stratí~--- a s~ňou svoje skóre zmeny,
          teda práve tú veličinu, pre ktorú štúdia existuje;
    \item žiak, ktorý neexistuje, ticho získa riadok.
\end{itemize}

Obe sú neviditeľné, a tu je to horšie než obvykle: analytický plán pokladá
výpadok za náhodný podmienene (MAR) a dopĺňa ho imputáciou. Imputovať hodnotu
neexistujúcemu žiakovi a zahodiť príležitosť skutočného ako neúčasť znamená
skresliť výsledok v~smere, ktorý nikto nedokáže podpísať.

\paragraph{Riešenie má tri vrstvy a každá zachytáva to, čo predchádzajúca nevie.}

\begin{description}[leftmargin=0.6cm, nosep]
    \item[\textbf{Kontrolný znak}] podľa ISO~7064 MOD~37,36. Zachytí každú
          jednoznakovú zámenu a každé prehodenie susedných znakov, teda veľkú
          väčšinu preklepov~--- a zachytí ich \emph{pri klávesnici}, kým je žiak
          ešte v~triede a môže kód prepísať.
    \item[\textbf{Abeceda bez zámenitelných znakov}] O/0, I/1/L, S/5, Z/2, B/8 sa
          nikdy nevydávajú. Znak mimo vydávanej abecedy teda nie je kódom vôbec a
          nástroj môže povedať, \emph{ktorá} chyba sa stala („v~kódoch nie je
          písmeno O~--- nemyslel si nulu?“), namiesto aby odmietol bez
          vysvetlenia. Trieda chýb sa tým mení na pokyn.
    \item[\textbf{Jeden kód, jedno sedenie}] kód už otvorený na inom počítači sa
          odmietne. Toto je jediná vrstva, ktorá zachytí chybu, ktorú kontrolný
          znak preukazateľne zachytiť nemôže: žiaka, ktorý zadá \emph{platný} kód
          niekoho iného. Aritmetika dva legitímne kódy od seba odlíšiť nedokáže.
\end{description}

\paragraph{Vlastnosť normy, ktorá nie je pravdivá.} ISO~7064 MOD~37,36 sa bežne
opisuje tak, že zachytí každú zámenu jedného znaku \emph{a} každé prehodenie
susedných znakov. Prvá polovica platí presne. Druhá nie: pri vyčerpávajúcom
preverení prejde približne jedno prehodenie z~dvoch tisíc, napríklad
\texttt{GY-YEDET-R} $\rightarrow$ \texttt{GY-YDEET-R}, pričom obidva kódy
kontrolou prejdú. Príčinou je pretečenie v~hybridnej rekurencii, ktoré pre jednu
dvojicu pozícií zopakuje účinný násobiteľ.

Na výber boli dve možnosti: vykázať 99{,}95\,\%, alebo urobiť z~toho 100\,\%. Je
to 100\,\%, pretože kódový priestor jednej školy má $9\,765\,625$ prvkov a škola
potrebuje asi tridsať: pri vydávaní sa vyskúšajú všetky prehodenia susedných aj
vzdialenejších znakov kandidáta a kandidát sa zahodí, ak hoci jedno z~nich
kontrolou prejde. Vlastnosť potom platí konštrukciou pre každý kód, ktorý sa
skutočne dostane do obálky, a test ju vyčerpávajúco preveruje namiesto toho, aby
veril norme na slovo. Zahodí sa asi jeden kandidát zo šiestich.

\paragraph{Čo to znamená pre štúdiu.} Pri pesimistickom predpoklade jedného
percenta chyby na znak má asi 5{,}9\,\% zadaní nejakú chybu, z~toho prevažná
väčšina jednoznakovú~--- a tie sa zachytia všetky. Očakávaný počet tichých
spojení dvoch žiakov za celú štúdiu je rádu $10^{-6}$. Bez kontrolného znaku by
nepozorovane prešlo približne sto zadaní.

\paragraph{Čo sa zámerne nerobí.} Žiadne samogenerované kódy (iniciály, mesiac
narodenia matky a podobne). Sú bežnou odpoveďou tam, kde zoznam neexistuje, a sú
súčasne slabšie a horšie: vybavenie po troch týždňoch je nedokonalé, takže
príležitosti sa stratia aj tak, a zložky kódu sú kvázi\-identifikátory, takže
„anonymný“ kód dokáže spätne priradiť ktokoľvek, kto má zoznam triedy. Školou
držaný zoznam je lepší v~oboch ohľadoch a táto štúdia ho má.

\section{Oprava kotvenia a to, čo by chyba stála}
\label{sec:kotvenie-oprava}

Metodika tvrdí, že \emph{šesť položiek je v~oboch formách totožných}, a že práve
tieto kotviace položky robia obe formy porovnateľnými. Príloha to tvrdí tiež:
kotviace položky sa vyskytujú „v~nezmenenej podobe“ aj v~okamžitom teste.

Ani jedna zo šiestich totožná nebola. Okamžitá forma sa pýtala na pole s~1024
prvkami, odložená na pole s~512; okamžitá uvádzala desaťprvkové pole, odložená
osemprvkové; jedna „kotviaca“ položka mala v~jednej forme štyri možnosti a v~druhej
voľnú numerickú odpoveď; dve boli rubrikové a ich zadanie bolo preformulované.
Kotviaca položka s~iným $n$ je \emph{iná} položka s~inou obtiažnosťou. Prevod
foriem na spoločnú škálu teda neexistoval, a rozdiel medzi okamžitým a odloženým
skóre~--- primárny kontrast práce~--- by zamieňal zabúdanie s~obtiažnosťou formy.

Ide o~tretí výskyt tej istej systémovej chyby: oprava bola urobená v~jednej
reprezentácii a nie v~druhej. Simulácia odhalila, že formy nie sú kotvené,
a oprava šla do \texttt{pipeline.py}; položky, ktoré by sa tým kódom
spracovávali, zostali nekotviteľné.

\paragraph{Koľko to stálo, číselne.} Pri vloženom zabúdaní \asmEqPlanted{}~logitu
a skutočne totožných kotviacich položkách sa hodnota obnoví na \asmEqEstZero{}
(skreslenie \asmEqRepairedBias). Ak sa šesť „kotviacich“ položiek líši a sú
o~pol logitu ľahšie~--- čo je pri 1024 oproti 512 striedmy odhad~--- obnoví sa
na \asmEqEstHalf{}.

\begin{table}[H]
\centering
\caption{Obnovenie vloženého zabúdania podľa rozdielu kotviacich položiek}
\label{tab:kotvenie-ucinok}
\small
\begin{tabular}{@{}lccl@{}}
\toprule
\textbf{Rozdiel kotviacich položiek} & \textbf{Odhad} & \textbf{Skreslenie} &
\textbf{Čo by práca vykázala} \\
\midrule
žiadny (opravený stav) & \asmEqEstZero  & \asmEqBiasZero  & zabúdanie, správne \\
0{,}3 logitu           & \asmEqEstThree & \asmEqBiasThree & zabúdanie, výrazne podhodnotené \\
0{,}5 logitu           & \asmEqEstHalf  & \asmEqBiasHalf  & \textbf{zlepšenie namiesto zabúdania} \\
0{,}8 logitu           & \asmEqEstEight & \asmEqBiasEight & zlepšenie, veľké \\
\bottomrule
\end{tabular}
\end{table}

Tretí riadok je ten, pre ktorý sa to vykazuje: \textbf{znak sa obráti}. Pri pol
logitu rozdielu by štúdia vykázala, že žiaci sa medzi okamžitým a odloženým
testom \emph{zlepšili}, hoci vložená pravda bola, že zabudli. A nič by na to
neupozornilo~--- model konverguje, zhoda položiek s~modelom je v~poriadku, interval
je úzky. Chyba, ktorá nevyzerá ako chyba, je v~metodologickej práci najhorší
možný druh.

\paragraph{Oprava a jej vedľajší zisk.} Obe formy majú teraz
\asmFormItems{}~položiek, jednu špecifikačnú tabuľku, rovnaký počet rubrikových
položiek (\asmOpenPerForm) a šesť znak po znaku totožných kotviacich položiek.
Skrátenie okamžitej formy zo šestnástich na \asmFormItems{} vynútil časový
rozpočet (časť~\ref{sec:cas-rozpocet}), ale metodicky je to zisk, nie ústupok:
izomorfia foriem je predpoklad, na ktorom primárny kontrast stojí, a dve formy
rôznej dĺžky s~rôznym pomerom formátov nie sú izomorfné v~žiadnom zmysle, ktorý
by čitateľ mohol overiť. Cenou je spoľahlivosť: dvanásť položiek dáva približne
$0{,}73$ namiesto $0{,}78$ pri šestnástich~--- rovnako ako odložená forma, ktorá
je primárnym výstupom.

Kotvia sa výhradne uzavreté položky. Pri rubrikovej položke by do prevodu foriem
na spoločnú škálu vstúpila aj zhoda medzi hodnotiteľmi, čiže chyba merania
hodnotenia by sa stala chybou prepojenia škál.

\section{Cena totožného kotvenia a rotovaný dizajn}
\label{sec:rotovane-kotvenie}

Totožné kotviace položky kotvenie opravia a prinesú vlastný problém, ktorý je
väčší, než sa dalo tušiť. Žiak tie šesť položiek videl pred tromi týždňami, takže
ich druhé zodpovedanie je \emph{vybavovacie precvičenie}. Kalibrácia nemá ako
vedieť, že položka je pre tohto žiaka ľahšia, a jej ľahkosť prečíta ako
schopnosť.

Pri precvičení 0{,}25~logitu sa vložené zabúdanie \asmEqPlanted{} odhadne na
\asmTeForgetting{}~--- takmer celé zmizne. Kontrast medzi podmienkami však
chránený \emph{je}: odhadne sa na \asmTeContrast{} pri vloženom
\asmTeContrastPlanted{} (skreslenie \asmTeContrastBias), pretože obe formy píše
každý žiak v~každej podmienke rovnako. Konfirmačná rodina práce je teda v~poriadku;
absolútna úroveň zadržania nie, a práca o~absolútnej úrovni tvrdenia vyslovuje.

\paragraph{Prečo pevná sada kotviacich položiek tento problém vyriešiť nemôže.}
Precvičenie by sa dalo odpočítať, keby sa poznala jeho veľkosť. Pri pevnej sade
ju poznať nemožno: každý žiak videl každú kotviacu položku, takže neexistuje
porovnávacia skupina. Problém nie je v~štatistike, ale v~dizajne.

\paragraph{Rotovaný dizajn.} Náhodné pridelenie porovnávaciu skupinu vyrobí. Téma
má zásobník \asmRotPool{}~kalibrovaných položiek, zloženého tak, aby každý
prípustný výber splnil špecifikačnú tabuľku (3~zapamätanie, 3~porozumenie,
6~aplikácia, 5~analýza, 1~hodnotenie). Pre každého žiaka sa náhodne vyberie šesť
položiek ako \emph{jeho} kotviaca sada (1/1/2/1/1 po úrovniach); okamžitá forma je
táto sada doplnená zo zásobníka na 2/2/4/3/1, odložená je tá istá sada plus
doplnok. Obe formy majú dvanásť položiek a obe splnia tabuľku, pre každého žiaka.

Dôsledok je ten, na ktorom záleží: v~odloženom teste odpovedá na každú položku aj
žiak, ktorý ju už videl, aj žiak, ktorý nie~--- a tie dve skupiny sú pridelením
ekvivalentné. Z~toho vyplývajú dve veličiny, ktoré pevný dizajn vyrobiť nedokáže:

\begin{description}[leftmargin=0.6cm, nosep]
    \item[\textbf{Veľkosť precvičenia}] je rozdiel logaritmov šancí medzi tými
          dvoma skupinami na tej istej položke. Odhadne sa na \asmRotPractice{}
          pri vloženom \asmRotPracticePlanted{}, a je \emph{identifikovaná
          pridelením}, nie predpokladom.
    \item[\textbf{Odložená schopnosť bez precvičenia}] sa odhadne len z~tých
          položiek, ktoré žiak predtým nevidel, pri obtiažnostiach pevne daných
          z~okamžitej príležitosti, kde precvičenie nie je žiadne. Je teda
          bezprecvičenia \emph{konštrukciou}, nie korekciou po fakte.
\end{description}

Zabúdanie potom vyjde \asmRotForgetting{} oproti vloženým \asmEqPlanted{}
(skreslenie \asmRotBias), kým pri pevnej sade vyjde \asmFixForgetting{}
(skreslenie \asmFixBias). Rotácia skreslenie zmenšuje približne trojnásobne a
naviac dodá odhad precvičenia samého.

\paragraph{Korekcia po fakte nefunguje, a stojí za to povedať prečo.} Prvý pokus
odhadol precvičenie zo náhodného kontrastu a pripočítal ho k~bežnému odhadu
zabúdania. Výsledok bol horší než neopravený, pretože bežný odhad už polovicu
precvičenia pohltil do presunu priemernej obtiažnosti~--- korekcia ju spočítala
druhýkrát. Druhý pokus odhadoval schopnosť, obtiažnosť a precvičenie spoločne
jednou vierohodnosťou a vrátil precvičenie $-378$~logitov: žiak, ktorý odpovedal
správne na všetky administrované položky, nemá konečný odhad schopnosti, Newtonov
krok pre neho beží na hranicu, a s~parametrom precvičenia v~tej istej
vierohodnosti ho ohraničení žiaci tlačia. Tretí pokus~--- obtiažnosti pevne
z~okamžitej príležitosti, schopnosť len z~nevidených položiek~--- funguje, pretože
nepotrebuje korekciu ani spoločnú vierohodnosť.

\paragraph{Cena a stav.} Zásobník \asmRotPool{}~položiek na tému už existuje:
toľko ich obe formy dohromady obsahujú. Čo neexistuje, je pilot~--- zásobník treba
najprv nakalibrovať, aby bol výber z~neho voči špecifikačnej tabuľke neutrálny~---
a na úrovni hodnotenia je v~zásobníku jediná položka, ktorú teda zdieľa každý
žiak, takže tam precvičenie odhadnuteľné nie je. \textbf{Ide o~odporúčanú zmenu
dizajnu s~pilotom, nie o~zmenu už vykonanú.} Primárnym dizajnom zostáva pevné
kotvenie s~priznaným obmedzením: kontrast medzi podmienkami je platný, absolútna
úroveň zadržania je nadhodnotená o~veľkosť precvičenia.

\section{Časový rozpočet počítaný z~položiek}
\label{sec:cas-rozpocet}

Práca mala časový rozpočet už predtým a vykazovala, že každé stretnutie sa
zmestí. Počítal každý blok ako \emph{počet položiek krát priemerný čas na
položku}, s~jedným priemerom na nástroj. To je v~poriadku, keď položky stoja
približne rovnako, a je to výrazne nesprávne, keď nie~--- a tu nie. Okamžitý test
mal desať uzavretých položiek po štyridsiatich sekundách a šesť rubrikových po
dvoch a pol minúte; plochý priemer 55~sekúnd rozdiel rozdelil a test podhodnotil
o~päť minút. Transferový test je šesť rubrikových položiek, ktorým príloha sama
priznáva dvadsaťpäť minút, a priemer im dal šesť a pol.

Chyby sa nevyrušia, pretože ukazujú rovnakým smerom: každý nástroj, ktorý priemer
podhodnocuje, je nástroj s~rubrikovými položkami, a tie sú zhromaždené na dvoch
miestach~--- v~odloženom teste, opakovanom pre tri témy, a v~transferovom teste.
Obe sedia v~stretnutí~T4, ktoré nesie primárny výstup štúdie.

Rozpočet sa preto počíta z~položiek, položku po položke. Celkové podhodnotenie
plochého priemeru je \asmBudgetUnderstated{}~minúty, najviac pri transferovom
teste o~\asmBudgetWorstBy{}~minúty. Pri takto počítanom rozpočte neprešli tri
z~piatich dozorovaných stretnutí.

\paragraph{Dve pravidlá, ktoré plochý priemer nemal ako vyjadriť.}

\begin{description}[leftmargin=0.6cm, nosep]
    \item[\textbf{Únavový strop}] meracie stretnutie nie je len súčet minút. Nad
          približne \instrFatigueCap{}~minút nepretržitého testovania výkon
          dospievajúcich klesá, a klesá rýchlejšie u~žiakov s~nižšou motiváciou~---
          čo je veličina, ktorú podmienky pravdepodobne ovplyvňujú. Priveľmi dlhé
          stretnutie teda nepridáva len šum, ale šum korelovaný so zásahom.
          Strop je požiadavka validity, nie ohľaduplnosti, a platí na najdlhší
          \emph{neprerušený} úsek, nie na súčet: tri pätnásťminútové testy
          s~prestávkami nie sú štyridsaťpäťminútový test.
    \item[\textbf{Náklad odovzdania}] každé dozorované stretnutie strávi čas,
          ktorý nie je testovaním: rozdanie kódov, pokyny, prebudenie notebookov.
          Účtuje sa raz za stretnutie a líši sa podľa stretnutia~--- T0a rozdáva
          obálky a vysvetľuje štúdiu žiakom, ktorí ju nikdy nevideli, kým pri P1
          už žiaci postup poznajú a majú sa len prihlásiť.
\end{description}

\paragraph{Čo z~toho vyplynulo pre dizajn.} Štyri zmeny, z~ktorých tri sú
metodickými zlepšeniami a jedna je náklad, ktorý nesie škola:

\begin{enumerate}[nosep]
    \item \textbf{Okamžitá forma má \asmFormItems{} položiek namiesto
          šestnástich.} Vedľajším ziskom je jedna špecifikačná tabuľka pre obe
          formy (časť~\ref{sec:kotvenie-oprava}).
    \item \textbf{Predtest k~téme sa presunul na konec predchádzajúcej hodiny
          informatiky} (4~min, pod dozorom, 2--3 dni vopred). Vedomosti
          o~predpokladoch témy sa za dva dni bez výučby nemenia, takže kovariát
          nič nestráca, a blok získa štyri minúty, ktoré v~ňom nie sú.
    \item \textbf{Výklad je 64~minút, nie 75.} Kontrolný list vernosti (V1--V3)
          sčítal 35+40+15 = 90~minút vyučovania vnútri 90-minútového bloku,
          ktorý má okrem toho obsahovať aj test a dotazník. Rozvrh a kontrolný
          list si protirečili o~pätnásť minút, a obe tvrdenia boli v~tej istej
          prílohe.
    \item \textbf{T4 sa delí na T4a a T4b.} Spolu je to 92~minút merania a
          transfer je práve to meranie, ktorému únava uškodí najviac: šesť
          rubrikových položiek vyžaduje odvodenie, nie vybavenie. Počet meracích
          hodín, ktoré škola dáva, stúpa na \asmMeasureLessons{}.
\end{enumerate}

Po týchto zmenách sa zmestia všetky stretnutia a žiadne neprekračuje únavový
strop. Rozpočet je v~tabuľke~\ref{tab:budget-gen}, generovanej zo zásobníka.

\paragraph{Na čom tento rozpočet stojí, a je to odhad.} Rozpočet je presnejší než
plochý priemer, pretože počíta položku po položke. Nie je však meraním: čas na
položku (\texttt{expected\_sec}) je \emph{úsudok} autora o~tom, koľko potrebuje
pozorný žiak~--- odvodený z~typu položky, dĺžky zadania a počtu krokov, ale
neoverený na žiakoch. Celé rozhodnutie o~rozdelení stretnutia~T4 na dve teda
stojí na číslach, ktoré ešte nikto nenameral.

Povedať to je potrebné, pretože inak by rozpočet vyzeral ako zistenie, hoci je
hypotéza. Dôsledky sú dva. Prvý: pilot tieto čísla nahradí pozorovanými mediánmi
z~jednej triedy, a rozpočet sa prepočíta~--- a keď sa prepočíta, môže sa ukázať,
že T4 sa deliť nemusí, alebo naopak, že sa musí deliť aj vyučovací blok. Druhý a
dôležitejší: aj keby boli všetky časy o~tretinu nadhodnotené, \emph{poradie}
nástrojov podľa náročnosti sa nezmení, a to je to, čo rozhodlo. Plochý priemer
pokazil práve poradie: postavil transferový test, najnáročnejší nástroj v~pláne,
pred predtest. Chyba v~úrovni je iná vec než chyba v~poradí, a opravená je zatiaľ
druhá.

Nástroj je na toto pripravený: \texttt{expected\_sec} je pole v~zásobníku a
rozpočet je funkcia nad ním, takže nahradenie odhadov meraniami je jedna úprava
zásobníka a jedno opätovné spustenie, nie prepísanie plánu.

\section{Skórovanie: čo robí stroj a čo nerobí}
\label{sec:skorovanie}

Uzavreté položky skóruje modul, okamžite a bez človeka. Rubrikové položky
\textbf{nie}, a odmietnutie je zámerné. Automatické hodnotenie voľného
slovenského textu o~algoritmickom uvažovaní by bolo nespoľahlivé spôsobom, ktorý
je horší než šum: bolo by nespoľahlivé \emph{diferenciálne}. Žiak, ktorý strávil
hodinu čítaním plynulej prózy asistenta, píše plynulejšiu prózu, a akýkoľvek
ukazovateľ povrchových príznakov by to odmenil. Hypotéza práce je, že podmienky
$A$ a $G$ sa líšia v~tom, \emph{ako} žiaci píšu. Hodnotič citlivý na to, ako
žiaci píšu, by teda časť účinku, ktorý meria, vyrábal~--- čo je presne tá chyba,
ktorú klasifikátor PEQI už raz mal a musel byť opravený.

Rubrikové položky teda odchádzajú ako \texttt{pending}, s~rubrikou pripojenou,
a vracajú sa dvojito kódované. Čo modul zaručuje, je úplnosť odovzdania: každá
voľná odpoveď sa exportuje na kódovanie, každé kódovanie sa načíta späť, a
\texttt{finalise()} odmietne vyrobiť analytický súbor, kým je hoci jedna
rubriková položka nekódovaná. Tichá nula za nekódovanú položku by vyzerala ako
nesprávna odpoveď.

\paragraph{Tri veci, ktoré skórovací skript obvykle nerobí.}

\begin{description}[leftmargin=0.6cm, nosep]
    \item[\textbf{Čiastočné body oddelene od dichotómie.}] Raschov model je
          dichotomický. Rubriková položka za 0--2 body sa preto exportuje aj ako
          hrubé body, aj ako dichotomizácia na strede rubriky, a ktorú z~nich
          analýza použije, je rozhodnutie analýzy a je deklarované tam. Čo
          prípustné nie je, je aby si ju tento modul ticho vybral sám.
    \item[\textbf{Zaznamenáva sa rozptyľovač, nie len nesprávnosť.}]
          \texttt{correct = 0} zahodí jedinú diagnostickú informáciu, ktorú
          položka s~voľbou vyrobí. Zvolená možnosť aj jej diagnóza putujú
          s~riadkom, a práve to umožňuje triednej správe povedať učiteľovi,
          \emph{ktorú} mylnú predstavu trieda drží. Zásobník nesie
          \asmDistractorsDiagnosed{}~diagnóz.
    \item[\textbf{Skórovanie je oddelené od vylúčenia.}] Odpoveď daná za
          1{,}4~sekundy sa skóruje normálne a označí sa samostatne
          (časť~\ref{sec:usilie}). Spojiť to dvoje by znamenalo skryť rozhodnutie
          o~vylúčení údajov do skórovacej funkcie, kde ho žiadny oponent hľadať
          nebude.
\end{description}

\paragraph{Nezhoda kódovateľov sa nepriemeruje.} Dve nezávislé hodnotenia rubriky,
ktoré sa rozchádzajú, sa riešia konsenzom. Priemer dvoch čítaní rubriky je číslo,
ktoré by nepodpísal ani jeden z~kódovateľov, a podiel nezhôd je sám vykazovanou
veličinou (váženou Cohenovou $\kappa$).

\section{Úsilie pri teste}
\label{sec:usilie}

Práca už pomenúva zámenu: „žiaci sa naučili viac“ a „žiaci sa na teste po hodine,
ktorá ich bavila, viac snažili“ sú ten istý kontrast, pokiaľ ich niečo neoddelí.
Jej dosiaľ jediná odpoveď je negatívna kontrola~--- výstup, ktorý sa nesmie
pohnúť. Je to správny druh odpovede a je neúplná, pretože negatívna kontrola
povie, \emph{že} je niečo zlé, nie \emph{čo}. Ak zlyhá, štúdii zostane
nevysvetlené zlyhanie a žiadny spôsob, ako rozhodnúť, či šlo o~úsilie, únavu
alebo pokazený nástroj.

Úsilie pri teste sa dá merať priamo a je to tu najpravdepodobnejší kandidát, a to
z~dôvodu špecifického pre tento dizajn: podmienka $G$ odstraňuje úsilie
\emph{počas hodiny} a predpokladá sa pri nej najvyššia spokojnosť. Žiak, ktorý mal
najpríjemnejšiu hodinu, potom píše test, ktorý sa neznámkuje. Či doň investuje, je
empirická otázka, a odpoveď sa pravdepodobne líši podľa podmienky~--- čím je
diferenciálna, a diferenciálne úsilie na výstupe nie je šum. Je to skreslenie
v~smere hypotézy alebo proti nej, a bez merania nikto nepovie, v~ktorom.

Meranie vyžaduje čas na položku, čo vyžaduje digitálnu administráciu. Toto je
najsilnejší jediný argument pre \texttt{runner.py}: papierový test toto číslo
nevyrobí vôbec.

\paragraph{Ako sa identifikuje rýchle hádanie.} Odpoveď rýchlejšia, než sa
položka dá prečítať, nie je odpoveďou. Hranica je miesto, kde sa metóda obvykle
ohýba, takže sú implementované tri pravidlá a vykazujú sa všetky tri: desať
percent očakávaného času položky (bežná konvencia), pevná dolná hranica (istota
proti položke s~priveľmi nízkym očakávaným časom), a dolina medzi dvoma vrcholmi
pozorovaného rozdelenia časov tam, kde nejaká je. Tretie pravidlo je jediné, ktoré
používa dáta, a vracia „žiadna hranica“, keď je rozdelenie jednovrcholové~---
metóda, ktorá by vždy vrátila číslo, by vyrobila rýchle hádanie v~triede, ktorá
žiadne nerobila. Nezhoda pravidiel je informácia: výsledok, ktorý platí pri
všetkých troch, je voči hranici odolný; výsledok, ktorý sa objaví pri jednom, je
výsledkom o~hranici.

\paragraph{Operačné vlastnosti.} Rozdiel v~úsilí 0{,}08 medzi podmienkami~--- taký,
ktorý by vedel predstierať najmenší zaujímavý účinok~--- sa odhalí
v~\asmEffortDetect{} prípadov, pri falošnej pozitivite \asmEffortFalsePos{} tam,
kde rozdiel nie je.

\paragraph{Čo sa nerobí.} Odpovede označené ako rýchle hádanie sa neodstraňujú;
označenie je samostatná premenná. Odstránenie by ticho zmenilo odhadovanú veličinu
z~„čo žiaci dosiahli“ na „čo dosiahli, keď sa snažili“, a tie dve sa líšia práve
vtedy, keď na tom záleží najviac. Vopred stanovené pravidlo: hlavná analýza na
všetkých odpovediach, analýza po filtrovaní úsilia ako deklarovaná robustnostná
kontrola, obe vždy.

\section{Kontaminácia počas testu}
\label{sec:kontaminacia}

Toto je štúdia o~tom, čo sa stane s~učením, keď má žiak asistenta s~umelou
inteligenciou. Jej výstupom je test. Test sa píše na notebooku, v~škole, žiakom,
ktorý má telefón. Každá publikovaná štúdia tohto druhu to rieši rovnako: dozorca
odfajkne položku. V~tejto práci je tou položkou V12 kontrolného listu vernosti~---
„počas testu žiaci nemali prístup k~asistentovi ani k~poznámkam“~--- a je to
jediné, čo stojí medzi primárnym výstupom a kontaminovaným.

A zlyháva v~nesprávnom smere. Žiak, ktorého hodina frustrovala, má menej dôvodov
starať sa o~neznámkovaný test než žiak, ktorého bavila; a žiak, ktorý hodinu
dostával hotové odpovede, má aj návyk, aj očakávanie, že si jednu vyhľadá.
Kontaminácia je teda pravdepodobne \emph{diferenciálna podľa podmienky}, čím je
skreslením, nie šumom, a skreslením, ktorého znak sa nedá uhádnuť.

\paragraph{Päť rodín príznakov z~troch nezávislých zdrojov.} Okno testu stratilo
zameranie; text odpovede bol vložený, nie napísaný; položka správne a výrazne
rýchlejšie, než sa dá prečítať; voľná odpoveď takmer totožná s~odpoveďou iného
žiaka; ťažšie položky žiaka trvali kratšie než ľahšie; a záchytná vrstva
zaznamenala dopyty na asistenta v~čase testu. Každý jednotlivo je slabý. Spolu
tiež nie sú silné~--- a práve preto je výstupom odstupňovaný príznak a nikdy
verdikt.

\paragraph{Tri pravidlá, ktoré z~toho robia použiteľnú vec a nie podozrievanie.}

\begin{enumerate}[nosep]
    \item \textbf{Detektor je kalibrovaný na stretnutí, kde kontaminácia nie je
          možná.} T0a prebehne skôr, než žiak nástroj videl, v~miestnosti, kde
          žiadny asistent nie je. Čokoľvek tam detektor označí, je falošná
          pozitivita \emph{konštrukciou}. Táto hodnota sa vykazuje, a detektor,
          ktorý označí 8\,\% na T0a, je detektor, ktorého 9\,\% na T4a nič
          neznamená. Samotný podiel príznakov v~poli nehovorí nič~--- je to tá
          istá chyba ako vykazovať podiel pozitívnych testov bez špecificity.
    \item \textbf{Žiadny žiak sa škole nepomenuje.} Ide o~meranie kvality údajov,
          nie o~obvinenia o~deťoch, a rozdiel nie je rečnícky: príznak, ktorý by
          sa mohol dostať ku známke, by zmenil motiváciu žiaka písať test
          poctivo, a tým by zničil meranie, pre ktoré bol zavedený. Správa pre
          školu nesie podiel za triedu a nič iné.
    \item \textbf{Vylúčenie je stanovené vopred a vykazuje sa vždy dvojmo.}
          Odpoveď sa vylúči len pri dvoch príznakoch z~rôznych rodín nad hranicou.
          Hlavná analýza beží na všetkom, analýza s~vylúčením vedľa nej, a obe sa
          vytlačia, akokoľvek vyjdú.
\end{enumerate}

\paragraph{Operačné vlastnosti.} Pri realistickom nazeraní~--- žiak si pozrie dve
až štyri položky z~dvanástich, nie všetky~--- detektor zachytí
\asmIntegDetect{} kontaminovaných žiakov v~triede
pri falošnej pozitivite \asmIntegFalsePos{}. To nie je dozor a nemá ním byť: je to
dolná hranica, ktorá z~nemerateľného tvrdenia robí vykazovateľný podiel.

\paragraph{Štyri rodiny v~deň zberu, piata až v~analýze.} Toto rozlíšenie sa
ukázalo pri skúšobnom prejdení celého školského priečinka a stojí za to ho
uviesť, pretože inak by sa od triednej správy čakalo viac, než vie dať. Rodiny
B,~T,~S a~C pracujú s~veličinami, ktoré sú známe okamžite~--- stav okna, čas
odpovede, podobnosť textov, záznam záchytnej vrstvy~--- a dajú sa teda vyhodnotiť
v~ten istý deň. Rodina~I porovnáva čas odpovede s~\emph{obtiažnosťou} položky, a
obtiažnosť pred Raschovou kalibráciou neexistuje. Beží preto až v~analytickej
vrstve, nie v~školskom priečinku. Správa, ktorú škola dostane v~deň zberu, je
teda založená na štyroch rodinách z~piatich, a vykazuje to.

Prvá verzia tejto simulácie dávala kontaminovanému žiakovi odpovedať na každú
položku za 3{,}5~sekundy s~osemnástimi sekundami straty zamerania, a detektor
proti nej dosiahol 100\,\%. To nehovorí nič, pretože nič také zjavné sa
v~miestnosti s~učiteľom nedeje. Lichotivé číslo je horšie než striedme.

\section{Kategorizácia žiakov, ktorá funguje pre jednu triedu}
\label{sec:pravidlovy-profil}

Zadanie žiada od nástroja kategorizáciu používateľov. Metodika odpovedá latentnou
profilovou analýzou nad šiestimi ukazovateľmi, a odpovedá dobre: profily sú
pomenované vopred, kovariančná štruktúra bola vybraná oproti vloženej pravde aj
oproti vloženej neprítomnosti štruktúry, a typológia sa poctivo vykazuje ako
exploratívna.

Len si treba všimnúť, čo to znamená pre toho, kto by nástroj skutočne používal.
Latentná profilová analýza je tvrdenie o~\emph{vzorke}. Potrebuje stovky žiakov,
prikladá zmes tomu rozptylu, aký vzorka náhodou má, a označenie jedného žiaka
závisí od toho, kto je v~súbore s~ním. Učiteľ s~jednou triedou dvadsiatich dvoch
ju spustiť nedokáže, a keby ju spustili dve školy, dostali by dve rôzne a
neporovnateľné typológie.

Práca teda dodáva kategorizáciu, ktorú vie použiť len štúdia. Pre účel, ktorý
zadanie samo pomenúva~--- nástroj nasaditeľný v~škole~--- je to nesprávny objekt.

\paragraph{Pravidlo s~absolútnymi hranicami.} Nástroj preto obsahuje
deterministické, vopred stanovené pravidlo nad \emph{absolútnymi} hranicami. Štyri
dôsledky, z~ktorých žiadny zmesový model mať nemôže: funguje pri $n=1$; označenie
nezávisí od ostatných v~triede, takže dve školy v~rôznych rokoch dostanú
porovnateľné označenia; každá hranica je napísaná skôr, než existujú dáta, takže
nezostáva voľnosť ladiť typológiu, kým nepovie niečo zaujímavé; a učiteľovi sa dá
jednou vetou povedať, \emph{prečo} bol žiak označený, čo posteriórna
pravdepodobnosť zo šesťrozmernej gaussovskej zmesi nedokáže.

Cena, povedaná priamo: pravidlo nedokáže objaviť profil, ktorý nikto nečakal, a
zmes dokáže. Nie sú to konkurenti. Zmes zostáva v~práci ako výskumná otázka „je
táto populácia zmesou, a akou“; pravidlo je to, čo dostane škola.

\paragraph{Odkiaľ sa berú hranice.} Nie z~dát a nie z~konvencie. Každá je bodom na
škále, ktorá už má význam. PEQI je priemer úrovní promptov s~váhami
$L_0=0$ až $L_4=1$, takže $0{,}25$ je žiak, ktorého priemerný prompt je holá
žiadosť o~overenie, a $0{,}50$ ten, ktorý sa v~priemere pýta na vysvetlenie~---
obe hranice sú položené práve tam, pretože tam sa mení konštrukt, nie tam, kde sa
delí vzorka. Podiel úloh s~vlastným pokusom má hranicu $0{,}50$, čo je jediný
nearbitrárny bod na podiele. Kalibračné skreslenie má hranicu $0{,}15$, čo je
najmenšie nadhodnotenie prevyšujúce chybu merania dvanásťpoložkového testu.

\paragraph{Pravidlo, ktoré odmietne odpovedať.} Ukazovateľ pod svojím minimom
dôkazu je \emph{neprítomný}, nie len poznamenaný. Prvá verzia ho uvádzala ako
chýbajúci a potom ho aj tak použila, a nesprávne označila práve toho žiaka, pre
ktorého to rozlíšenie existuje: žiak s~jediným promptom mal PEQI $0{,}50$
spočítané z~toho jedného promptu, prešiel hranicou „pýta sa na vysvetlenie“ a bol
označený ako kritický používateľ na základe jednej otázky. Pravidlo, ktoré vždy
vydá označenie, je pravidlo, ktoré háda, a hádanie vytlačené pri mene žiaka je
horšie než prázdne miesto.

Podmienka $K$ nemá asistenta, takže profil používania nástroja preň nie je
definovaný. To nie sú chýbajúce údaje, ale podmienka, v~ktorej konštrukt
neexistuje.

\paragraph{Čo pravidlo dokáže a čo nie, číselne.}

\begin{table}[H]
\centering
\caption{Pravidlo a latentná profilová analýza oproti vloženej pravde (upravený Randov index)}
\label{tab:pravidlo-vs-lpa}
\small
\begin{tabular}{@{}lccc@{}}
\toprule
\textbf{Počet žiakov} & \textbf{Pravidlo} & \textbf{Zmesový model} &
\textbf{Počet komponentov} \\
\midrule
\asmNClass{} (jedna trieda) & \asmRuleAriClass & \asmLpaAriClass & \asmLpaKClass \\
\asmNThree{}                & \asmRuleAriThree & \asmLpaAriThree & \asmLpaKThree \\
\asmNStudy{} (táto štúdia)  & \asmRuleAriStudy & \asmLpaAriStudy & \asmLpaKStudy \\
\asmNLarge{}                & \asmRuleAriLarge & \asmLpaAriLarge & \asmLpaKLarge \\
\bottomrule
\end{tabular}
\end{table}

Pri jednej triede zmesový model nenájde žiadne profily~--- vyberie jediný
komponent~--- kým pravidlo dosiahne \asmRuleAriClass. Pri \asmNStudy{} a viac je
zmes lepšia. To je presne to rozdelenie úloh, ktoré práca navrhuje.

\emph{Pozor na porovnateľnosť:} tieto hodnoty nie sú porovnateľné s~upraveným
Randovým indexom, ktorý uvádza časť~\ref{sec:kategorizacia}. Tam sú ukazovatele
v~štandardných jednotkách s~vnútroprofilovým rozptylom $0{,}62$; tu sú na
absolútnych škálach, ktoré pravidlo potrebuje, a s~menším vnútroprofilovým
rozptylom. Porovnávať sa dajú obidva stĺpce tejto tabuľky medzi sebou, nie
s~hodnotou z~kapitoly~\ref{chapter:metodika}.

\paragraph{Označenie nezávisí od vzorky.} Ten istý žiak preklasifikovaný
v~\asmLabelCohorts{} rôznych náhodných kohortách dostal vždy to isté označenie:
\asmLabelStable{}. Zmesový model takú vlastnosť nemá a mať nemôže.

\paragraph{Citlivosť na hranice, vykázaná aj tam, kde nelichotí.} Posun jednej
hranice o~$\pm$20\,\% mení zhodu najviac o~\asmCutSingle{} (najcitlivejšia je
hranica \asmCutWorst), takže nezhoda o~jednu hranicu typológiu nerozvráti. Posun
všetkých piatich naraz a v~tom istom smere ju mení o~\asmCutJoint{}, a to málo
nie je: pravidlo s~absolútnymi hranicami \emph{je} citlivé na ich spoločné
nastavenie. Práve preto musia byť hranice odvodené z~významu škály a stanovené
vopred. Predregistrácia tu nie je formalita, ale nosný prvok~--- bez nej by sa päť
hraníc dalo ladiť, kým typológia nepovie to, čo sa od nej čaká.

\paragraph{Pridáva profil niečo nad vstupnú úroveň?} Ak je „slepo dôverujúci“ len
inak pomenované „nízke vstupné skóre“, typológia nepridáva nič. Test oproti
dvom vloženým svetom~--- v~jednom má profil vlastný účinok, v~druhom je presne
preznačkovanou vstupnou úrovňou~--- má silu \asmIncPower{} nájsť vlastný účinok,
keď tam je, a nájde ho v~\asmIncFalsePos{} prípadov, keď tam nie je. Bez tohto
testu by typológia mohla byť preznačkovaná vstupná úroveň a z~výsledku by to nikto
nerozlíšil.

\section{Čo vychádza von a komu}
\label{sec:spatna-vazba}

Protokol škole sľubuje „agregovanú správu za triedu (nie za žiaka) s~porovnaním
voči ostatným zapojeným triedam“. Čo v~nej je, nebolo špecifikované. To nie je
kozmetické: škola, ktorá nedostane nič užitočné, nepustí štúdiu na ďalší rok, a
merací nástroj, ktorý tým, koho meria, nevracia nič, je extraktívny. Zadanie žiada
nástroj, ktorý vedia používať školy; nástroj, ktorého jediným výstupom je CSV pre
výskumníka, ním nie je.

Ostrejší dôvod je ale iný. Triedna správa je \emph{zverejnenie}. Dvadsaťdva
žiakov, tri podmienky, štyri profily~--- tabuľka počtov krížená podmienkou má
bunky s~dvoma a tromi žiakmi, a učiteľ, ktorý vie, kto sedel v~ktorej podmienke,
z~malej bunky často prečíta jednotlivca. Potrebné sú teda pravidlá potlačenia,
ktoré uplatňuje kód a nie dobrý úmysel, a ktoré sa dajú otestovať.

\begin{description}[leftmargin=0.6cm, nosep]
    \item[\textbf{Učiteľ}] dostane, čo učiť na ďalšej hodine: počty mylných
          predstáv po položkách, na úrovni triedy, \emph{bez rozpisu podľa
          podmienky}. Podmienka je premenná výskumníka, učiteľovi pedagogicky
          neslúži, a jej pridanie je práve to, čo vyrobí malé bunky. Zoradené je
          to podľa toho, koľko žiakov drží tú istú mylnú predstavu, nie podľa
          obtiažnosti: položka, ktorú devätnásť žiakov pokazilo devätnástimi
          rôznymi spôsobmi, je zlá položka, kým položka, ktorú jedenásť žiakov
          pokazilo rovnako, je učivo.
    \item[\textbf{Škola}] dostane niekoľko agregátov v~porovnaní s~ostatnými
          triedami, s~potlačením buniek pod piatimi žiakmi, bez čohokoľvek na
          úrovni žiaka. Trieda pod dvanásť žiakov sa s~kohortou neporovnáva vôbec:
          pri malej triede je priemer triedy tvorený niekoľkými konkrétnymi deťmi
          a „vaša trieda je pod priemerom“ je potom tvrdenie o~troch z~nich.
          Neporovnáva sa ani proti kohorte menšej než tri triedy.
    \item[\textbf{Žiak}] nedostáva individuálne skóre, a je to rozhodnutie
          s~dôvodom: nástroje tejto štúdie majú spoľahlivosť vhodnú na porovnanie
          skupín, nie na výrok o~jednotlivcovi. Oznámiť žiakovi číslo, ktoré
          nástroj neunesie, by bola tá istá chyba, akú práca kritizuje inde.
          Po skončení štúdie dostáva prístup k~celej sade úloh a správnym
          odpovediam~--- informáciu, ktorú meranie unesie.
\end{description}

Pravidlá potlačenia sú overené \emph{auditom zverejnenia}, ktorý beží nad
skutočnými správami a nie nad zámerom: prejde celú štruktúru, hľadá početnosť pod
hranicou, čokoľvek, čo vyzerá ako pseudonym, a akékoľvek kríženie s~podmienkou.

\section{Čo musí rozhodnúť pilot}
\label{sec:modul-pilot}

Pilot na jednej triede je už súčasťou plánu (časť~\ref{sec:pilot}) a jeho
rozhodovacie pravidlá sú stanovené vopred. Hodnotiaci modul pridáva šesť otázok,
ktoré tam dosiaľ neboli, a pri každej aj to, čo jedna trieda rozhodnúť
\emph{nevie}.

\begin{table}[H]
\centering
\caption{Čo k pilotu pridáva hodnotiaci modul}
\label{tab:pilot-modul}
\scriptsize
\begin{tabularx}{\textwidth}{@{}X X l@{}}
\toprule
\textbf{Otázka} & \textbf{Vopred stanovené pravidlo} & \textbf{Rozhodne?} \\
\midrule
Sedia časy na položku?
& pozorované mediány nahradia odhady v~zásobníku a rozpočet sa prepočíta; ak sa niektoré stretnutie prestane zmestiť, delí sa blok s~najväčšou rezervou, nikdy nie primárny výsledok
& \textbf{áno} \\
Zvládnu pätnásťroční kód z~obálky?
& podiel zadaní odmietnutých kontrolným znakom sa vykazuje; nad 15\,\% sa kód skracuje na štyri znaky a cena v~kolíziách sa prepočíta
& \textbf{áno} \\
Dávajú rubriky zhodu medzi kódovateľmi?
& vážená $\kappa$ pod $0{,}67$ $\Rightarrow$ rubrika sa zlučuje na dvojúrovňovú; pod $0{,}50$ sa položka vyradí
& čiastočne \\
Aká je falošná pozitivita detektora kontaminácie na T0a?
& nad $0{,}10$ sa detektor z~vykazovania vypúšťa a uvedie sa dôvod
& \textbf{áno} \\
Dáva pravidlový profilovač použiteľné rozdelenie?
& podiel \emph{neurčený} nad $0{,}40$ $\Rightarrow$ minimá dôkazu sa prehodnotia \emph{pred} hlavným zberom a zmena sa zapíše
& slabo \\
Je zásobník \asmRotPool{}~položiek nakalibrovaný na rotáciu?
& bez kalibrácie rotované kotvenie do plánu nevstupuje
& \textbf{nie} \\
\bottomrule
\end{tabularx}
\end{table}

\paragraph{Prvý riadok je ten, na ktorom záleží najviac}, pretože na odhadoch
času stojí rozdelenie stretnutia~T4 (časť~\ref{sec:cas-rozpocet}). Pilot ich
premení z~úsudku na meranie, a je to jediná otázka z~tabuľky, ktorú jedna trieda
rozhodne úplne: medián času na položku pri dvadsiatich žiakoch je dosť presný na
rozhodnutie o~minútach, aj keď nie na nič iné.

\paragraph{Posledný riadok je ten, ktorý pilot rozhodnúť nevie}, a uvádzame ho
práve preto. Kalibrácia zásobníka na rotované kotvenie vyžaduje, aby každú
položku riešil dostatočný počet žiakov~--- pri osemnástich položkách a jednej
triede je to asi desať odpovedí na položku, čo nestačí ani na hrubú obtiažnosť.
Rotácia preto zostáva odporúčaním aj po pilote, pokiaľ sa nenájde druhá trieda
mimo hlavnej vzorky.

\paragraph{Čo pilot zmeniť nesmie} je rovnaký zoznam ako v~časti~\ref{sec:pilot},
rozšírený o~tri veci z~tohto modulu: hranice pravidlového profilovača, pravidlo
„dve rodiny z~rôznych rodín“ pre vylúčenie, a zásadu, že rubrikové položky
neskóruje stroj. Pilot slúži na nastavenie toho, čo sa dá nameraním zistiť, nie
na posúvanie rozhodnutí, ktoré boli stanovené vopred práve preto, aby sa
neposúvali.

\section{Čo tento modul nerieši}
\label{sec:modul-limity}

\begin{itemize}[nosep]
    \item \textbf{Správnosť rubrikového hodnotenia.} Modul zaručí, že každá
          voľná odpoveď bude dvojito kódovaná a že nezhoda sa vyrieši konsenzom.
          Nezaručí, že rubrika meria to, čo tvrdí.
    \item \textbf{Kontamináciu cez telefón.} Rodiny B, T a I vidia, čo sa deje
          v~okne testu a v~časoch odpovedí. Žiaka, ktorý má telefón pod stolom a
          odpíše pomaly, nevidí žiadna z~nich. Detektor je dolná hranica.
    \item \textbf{Validitu testu kritického myslenia.} Príloha~\ref{chapter:kriticke-myslenie}
          je návrh nástroja, nie validovaný nástroj, a pilot ešte neprebehol.
    \item \textbf{Absolútnu úroveň zadržania pri pevnom kotvení.} Je nadhodnotená
          o~veľkosť vybavovacieho precvičenia; opravuje to rotovaný dizajn, ktorý
          čaká na pilot (časť~\ref{sec:rotovane-kotvenie}).
    \item \textbf{Nič z~toho, čo sa stane, keď škola kód stratí.} Obnovenie
          zo~zoznamu robí kontaktná osoba školy; modul to nerobí a robiť nemá,
          pretože by na to potreboval presne to spojenie, ktoré nesmie existovať.
    \item \textbf{Rodinu~I v~deň zberu.} Porovnanie času s~obtiažnosťou položky
          vyžaduje kalibráciu, ktorá v~ten deň ešte neexistuje
          (časť~\ref{sec:kontaminacia}).
\end{itemize}

\section{Nasaditeľnosť: čo dostane škola a ako sa to overuje}
\label{sec:nasaditelnost-modul}

Zadanie žiada overenie riešenia \emph{použitím} v~školskom prostredí. Škola
nedostane repozitár, ale priečinok, a čo v~tom priečinku chýba, chýba v~ten deň.
Hodnotiaci modul je preto súčasťou školského priečinka (\texttt{build-school-kit.ps1}),
nie len repozitára:

\begin{description}[leftmargin=0.6cm, nosep]
    \item[\texttt{KODY.cmd}] vydá kódy žiakom. \textbf{Spúšťa ho škola, nie
    výskumník}, a vzniknú dva súbory: samotné kódy, ktoré idú výskumníkovi, a
    zoznam s~miestom pre meno, ktorý \emph{zostáva v~škole}. Keby kódy vydával
    výskumník, prepojovacia tabuľka by vznikla na jeho notebooku, a žiadna
    formulácia v~protokole by to nenapravila.
    \item[\texttt{TEST.cmd}] administruje test. Je to zámerne \emph{iný} spúšťač
    než \texttt{RUN.cmd} pre hodinu: hodina a test sú rôzne udalosti s~rôznymi
    pravidlami~--- počas testu musí byť asistent nedostupný~--- a jediný vstupný
    bod, ktorý by vedel spustiť oboje, je jeden chybný stisk od neplatného
    merania.
    \item[\texttt{SKONTROLUJ.cmd}] overí zozbierané súbory v~ten istý deň.
\end{description}

\paragraph{Priečinok sa overuje tým, že sa spustí, nie tým, že sa vymenuje.}
Zoznam súborov nezaručí, že sa priečinok dá použiť: moduly skopírované do plošej
štruktúry, než v~akej boli napísané, sa nemusia dať naimportovať, a to sa prejaví
až v~triede. Zostavenie priečinka preto na konci z~neho spustí Python, naimportuje
všetky moduly, načíta zásobníky, vydá kód a overí, že preklep v~ňom sa odmietne~---
a zostavenie zlyhá, ak hoci jedno z~toho neprejde.

\paragraph{Skúšobné prejdenie.} Celý priečinok bol prejdený tak, ako ho prejde
trieda: škola vydala kódy, šesť žiakov odpovedalo na všetkých
\asmFormItems{}~položiek okamžitého testu cez skutočný spúšťač, a výsledok sa
načítal do skórovača bez prekladovej vrstvy. Uzavreté položky boli oskórované,
rubrikové odišli na kódovanie, \texttt{finalise()} odmietol vytvoriť analytický
súbor, a triedna správa o~integrite vznikla bez jediného mena. Tento prechod
odhalil obmedzenie uvedené v~časti~\ref{sec:kontaminacia}~--- rodina~I v~deň zberu
nie je dostupná~--- ktoré by sa z~jednotkových testov nedalo vidieť, pretože tam
obtiažnosti položiek dodáva fixtúra.
